\documentclass[10pt,aspectratio=169]{beamer}

\usetheme{metropolis}
\usepackage{booktabs}
\usepackage{xcolor}
\usepackage{hyperref}

\definecolor{triggerblue}{HTML}{1F6FB4}
\definecolor{goodgreen}{HTML}{2E8B57}

\title{Macro News and a Futures Mapping Pipeline}
\subtitle{A review of capabilities, reasoning, and limitations}
% \author{James Zhang}
% \institute{}
\date{\today}

\begin{document}

\maketitle

% ---------------------------------------------------------------
\begin{frame}{Outline}
  \tableofcontents
\end{frame}

% ---------------------------------------------------------------
\section{Experiment background}

\begin{frame}{Experiment background}
  \scriptsize
  \begin{columns}[T,onlytextwidth]
    \begin{column}{0.50\textwidth}
      \textbf{Dataset}
      \begin{itemize}\itemsep0.1em
        \item Dow Jones Newswires, March 2022 (shards \texttt{a}--\texttt{e})
        \item 1{,}000 articles, random sample with fixed seed
        \item Month chosen for dense macro signal: Russia--Ukraine, Fed liftoff, commodity shock
      \end{itemize}
      \vspace{0.3em}
      \textbf{Asset universe}
      \begin{itemize}\itemsep0.1em
        \item 95 futures contracts across 7 asset classes
        \item Commodities, FX, government bonds, STIR, equity indices, US equity sectors, volatility
      \end{itemize}
      \vspace{0.3em}
      \textbf{Model \& infra}
      \begin{itemize}\itemsep0.1em
        \item Gemma 4 26B A4B Instruct via vLLM 0.19
        \item Deterministic ($T{=}0$), \texttt{max\_model\_len}${=}65{,}536$
        \item 1$\times$ RTX Pro 6000 Blackwell, $\sim$25 min end-to-end
      \end{itemize}
    \end{column}
    \begin{column}{0.48\textwidth}
      \textbf{Four-stage LLM pipeline}
      \begin{enumerate}\itemsep0.1em
        \item \textbf{Summarize} --- macro summary + company-specific gate
        \item \textbf{ArticleMapper} --- whole article $\rightarrow$ (asset, signal, score)
        \item \textbf{ParagraphMapper} --- each paragraph, with macro summary as context
        \item \textbf{Validator} --- reviews union, accepts/rejects each pair
      \end{enumerate}
      \vspace{0.4em}
      \textbf{Run totals}
      \begin{itemize}\itemsep0.1em
        \item 726 / 1000 filtered as company-specific (73\%)
        \item 4{,}157 (article, asset) candidates evaluated
        \item 4{,}019 accepted / 138 rejected (3.3\% reject)
        \item Signal mix: 727 strong / 3{,}292 weak
      \end{itemize}
    \end{column}
  \end{columns}
\end{frame}

% ---------------------------------------------------------------
\section{Example mappings}

\begin{frame}[fragile]{Example 1 --- BOE rate decision}
  \scriptsize
  \textbf{Headline:} \emph{BOE's Rate Decision Could Benefit Sterling in Short-Term --- Market Talk} \hfill \scriptsize(DJNW, 2022-03-16, shard \texttt{c})

  \vspace{0.4em}
  \begin{block}{Full article text (single ``Market Talk'' blurb)}
    \scriptsize
    1119 GMT --- Sentiment around sterling looks quite downbeat but the Bank of England's meeting
    on Thursday may offer the currency some brief relief, HSBC says. A 25 basis-point rate rise
    is fully priced so that's unlikely to have a notable impact on sterling but the currency will
    rise in the short-term if the split of votes on the decision shows some policymakers backing
    a larger increase, HSBC forex analyst Dominic Bunning says in a note. ``Should we see further
    calls for 50bp hikes this week by some Monetary Policy Committee members, then we would
    expect GBP to bounce.'' The BOE raised its key rate by 25 basis points to 0.50\% in February
    but four out of nine officials voted for a 50 basis-point rise. (renae.dyer@wsj.com)
  \end{block}

  \vspace{0.3em}
  \begin{columns}[T,onlytextwidth]
    \begin{column}{0.32\textwidth}
      \textbf{\textcolor{triggerblue}{Sterling}} \\[0.1em]
      [currency]\\
      strong, 0.95, both\\[0.2em]
      \emph{Transmission:} hawkish MPC split $\rightarrow$ higher expected
      terminal rate $\rightarrow$ wider rate differential $\rightarrow$ GBP up.
    \end{column}
    \begin{column}{0.32\textwidth}
      \textbf{\textcolor{triggerblue}{3-Month SONIA}} \\[0.1em]
      [STIR]\\
      strong, 0.95, both\\[0.2em]
      \emph{Transmission:} 50bp-vote signal directly repriced into the SONIA
      curve; SONIA is the sterling overnight benchmark.
    \end{column}
    \begin{column}{0.32\textwidth}
      \textbf{\textcolor{triggerblue}{UK Long Gilt}} \\[0.1em]
      [government bond]\\
      strong, 0.75, both\\[0.2em]
      \emph{Transmission:} more hawkish BoE path lifts gilt yields on the
      long end via the expectations and term-premium components.
    \end{column}
  \end{columns}

  \vspace{0.4em}
  \textcolor{goodgreen}{\textbf{Why this is a clean showcase:}} one trigger, three asset classes,
  both mappers independently agree, validator accepts all three. The pipeline correctly
  scores the gilt lower --- a real but less direct channel than FX/STIR.
\end{frame}

% ---------------------------------------------------------------
\begin{frame}[fragile]{Example 2a --- USDA soybean export data (article text)}
  \scriptsize
  \textbf{Headline:} \emph{U.S. Export Sales and Shipments: Soybeans --- Mar 17}
  \hfill (DJNW, 2022-03-17, shard \texttt{c})

  \vspace{0.3em}
  \textbf{[para 0, caption]} \emph{For the week ended Mar 10 in thousand metric tons.
  Source: USDA. The current marketing year for soybeans began Sep 1.}

  \vspace{0.2em}
  \textbf{[para 1, data table]} \emph{Totals may not add due to rounding.}

  \vspace{0.2em}
  \centering
  \setlength{\tabcolsep}{4pt}
  \renewcommand{\arraystretch}{1.0}
  \begin{tabular}{lrrrrrr}
    \toprule
     & \multicolumn{4}{c}{\textbf{2021--2022 (current MY)}} & \multicolumn{2}{c}{\textbf{2022--23}} \\
    \cmidrule(lr){2-5}\cmidrule(lr){6-7}
    \textbf{Country} & Wk $\Delta$ com. & Wk ship. & Seas. ship. & Seas. com. & Wk $\Delta$ com. & Seas. com. \\
    \midrule
    CHINA        & 395.5 & 370.2 & 25{,}626.3 & 27{,}686.3 & 406.0 & 5{,}639.0 \\
    EGYPT        & 241.8 &  61.8 &  2{,}316.7 &  3{,}349.7 &     0 &   105.0 \\
    TAIWAN       & 162.5 &  22.2 &     915.4 &  1{,}212.6 &     0 &       0 \\
    INDONESIA    &  83.2 &  87.7 &     922.3 &  1{,}201.3 &     0 &       0 \\
    MEXICO       &  56.2 & 120.6 &  2{,}812.9 &  4{,}805.6 &     0 &   123.5 \\
    \multicolumn{7}{l}{\emph{\ldots 40+ other countries, most 0}} \\
    \midrule
    \textbf{TOTAL KNOWN} & \textbf{985.6} & 714.3 & 42{,}325.5 & 49{,}456.2 & 406.0 & 6{,}115.8 \\
    \textbf{GRAND TOTAL} & 1{,}253.2 & 714.3 & 42{,}325.5 & 53{,}624.0 & 477.0 & 8{,}123.8 \\
    \bottomrule
  \end{tabular}

  \vspace{0.3em}
  \raggedright\footnotesize
  \emph{Fully structured data release --- no prose ``trigger'' sentence. The pipeline must read
  the table itself.}
\end{frame}

% ---------------------------------------------------------------
\begin{frame}{Example 2b --- reasoning \& transmission channel}
  \scriptsize
  \textbf{\textcolor{triggerblue}{Soybeans}} \quad [commodity] \quad \textbf{strong, 0.95, both}
  \quad \textbf{AM:} 0.95 \quad \textbf{PM(max):} 0.95 (para 1) \quad \textbf{evidence:} [1]

  \vspace{0.3em}
  \begin{block}{ArticleMapper reasoning}
    \scriptsize
    USDA weekly export shipment and commitment data is a primary fundamental driver for CBOT
    Soybean futures. It shows ``Week's net Shipments'' of \textbf{985.6k MT} and ``Season's
    Total Shipments'' of \textbf{42{,}325.5k MT}. It highlights significant activity in major
    demand centers like \textbf{China (395.5k MT)} and \textbf{Egypt (241.8k MT)}. These data
    points directly impact the global supply/demand balance and the pace of US soybean exports.
  \end{block}

  \begin{block}{ParagraphMapper [1] reasoning --- 0.95 (paragraph 0 only got 0.50)}
    \scriptsize
    The trigger is the ``total known shipments of 985.6k MT'' and the high-volume net export
    activity in China (395.5k MT) and Egypt (241.8k MT). Transmission: high weekly shipments
    indicate strong international demand and reduce immediate domestic supply, exerting upward
    pressure on soybean futures prices.
  \end{block}

  \begin{block}{Validator reasoning (accepted)}
    \scriptsize
    Direct report of USDA weekly data. The mappers correctly identified the China/Egypt figures
    and total known shipments as the primary triggers. \emph{``I agree with the ArticleMapper's
    assessment of a strong signal and high relevance.''}
  \end{block}

  \vspace{0.2em}
  \textcolor{goodgreen}{\textbf{Takeaways:}} (1) the pipeline reads a pure-data table and
  extracts the right country-level figures; (2) \textbf{evidence localization works} --- caption
  para 0 = 0.50, data para 1 = 0.95; (3) all three stages converge independently.
\end{frame}

% ---------------------------------------------------------------
\begin{frame}[fragile]{Example 3a --- Russia aerospace sanctions (article text)}
  \scriptsize
  \textbf{Headline:} \emph{Heard on the Street: Russia Can't Fly Without The West --- WSJ}
  \hfill (DJNW, 2022-03-09, shard \texttt{b})\\
  \textbf{Public link:} \href{https://www.wsj.com/business/airlines/russia-cant-fly-without-the-westbut-may-eventually-propel-china-11646738302}{\texttt{wsj.com/\ldots/russia-cant-fly-without-the-west\ldots}}

  \vspace{0.4em}
  \begin{block}{Relevant article paragraphs}
    \scriptsize
    \textbf{[3]} Russia operates 3\% of the global fleet, representing a steady stream of
    ``aftermarket'' cash for Western manufacturers. Boeing and Airbus have engineering centers
    there, which have now suspended operations. The U.S. company said Monday it had stopped
    buying from its Russian partner \textbf{VSMPO-AVISMA, the supplier of a quarter of the
    world's titanium} --- a key input for airframes.\\[0.3em]
    \textbf{[4]} Western lessors own \textbf{70\% of Russia's 981 jets} in service or in storage
    and have no clear way of repossessing them. But the hit is deeply asymmetrical.
    \textbf{Only 1\% of Airbus' and Boeing's order backlogs is affected.}\\[0.3em]
    \textbf{[7]} Yet both use a lot of Western parts and engines: The MC-21 flies geared
    turbofans built by Pratt \& Whitney. While work to ``russify'' the plane is already in the
    works, it would take decades\ldots\ China, which has invested heavily in a similar plane,
    the C919, is the obvious partner.
  \end{block}

  \vspace{0.3em}
  \emph{A rich, narrative WSJ column. Multiple macro angles: a commodity supply shock (titanium),
  aerospace-manufacturer revenue loss, and a long-horizon Russia--China technology pivot. The
  pipeline tags} \textbf{6 assets} \emph{and rejects a seventh --- we will walk through all of
  them.}
\end{frame}

% ---------------------------------------------------------------
\begin{frame}{Example 3b --- all 9 mappings the pipeline produced}
  \scriptsize
  \centering
  \setlength{\tabcolsep}{5pt}
  \renewcommand{\arraystretch}{1.12}
  \begin{tabular}{llcccl}
    \toprule
    \textbf{Asset} & \textbf{Class} & \textbf{Signal} & \textbf{Score} & \textbf{Source} & \textbf{Channel (pipeline's own words)} \\
    \midrule
    \textcolor{goodgreen}{\textbf{Materials Sector}} & US equity sector & \textbf{strong} & \textbf{0.75} & both & titanium supply shock (25\% world supply) \\
    Industrials Sector          & US equity sector & weak & 0.50 & both    & Boeing aerospace input costs \\
    Euro Stoxx 50 Index         & equity index     & weak & 0.45 & article & Airbus heavyweight in Eurostoxx \\
    Gold                        & commodity        & weak & 0.45 & para    & geopolitical risk premium \\
    Information Tech.\ Sector   & US equity sector & weak & 0.40 & article & Pratt \& Whitney / MC-21 supply chain \\
    CAC 40 Index                & equity index     & weak & 0.40 & article & Airbus (CAC 40 constituent) \\
    Mini Nikkei 225 Index       & equity index     & weak & 0.35 & article & Japanese aerospace component exposure \\
    DAX Index                   & equity index     & weak & 0.35 & article & German aerospace supply chain \\
    Hang Seng Index             & equity index     & weak & 0.35 & article & Russia--China bloc / C919 partnership \\
    \bottomrule
  \end{tabular}

  \vspace{0.4em}
  \textcolor{goodgreen}{\textbf{What to notice:}}
  \begin{itemize}\itemsep0.1em
    \item \textbf{Only one strong tag.} Materials Sector gets promoted because titanium is a
          quantifiable, article-specific driver. Everything else stays weak.
    \item Weak mappings follow the actual \textbf{economic neighborhood}: Boeing's Industrials,
          Airbus-heavy Euro indices (CAC 40, DAX, Euro Stoxx 50), Japanese aerospace via
          Nikkei, HSI via the C919 pivot.
    \item \textbf{No runaway tag.} The article caps the impact at 1\% of Boeing/Airbus order
          backlogs; the pipeline respects that --- every weak mapping sits in a tight
          0.35--0.50 score band.
  \end{itemize}
\end{frame}

% ---------------------------------------------------------------
\begin{frame}{Example 3c --- the Materials Sector reasoning, verbatim}
  \scriptsize
  \begin{block}{Materials Sector --- strong, 0.75, both (\textcolor{goodgreen}{accepted})}
    \scriptsize
    \textbf{ArticleMapper:} ``The article identifies a specific supply chain disruption involving
    a critical industrial input: titanium. Boeing stopped buying from its Russian partner
    VSMPO-AVISMA, the supplier of a quarter of the world's titanium --- a key input for
    airframes. This creates a direct supply shock to the global titanium market. Because
    titanium is a fundamental raw material for the aerospace and materials sectors, a sudden
    loss of 25\% of global supply is a significant trigger\ldots'' \\[0.3em]
    \textbf{ParagraphMapper [3]:} ``Boeing's decision to stop purchasing from VSMPO-AVISMA,
    a supplier responsible for 25\% of the world's titanium. The transmission channel is a
    direct supply-side shock to a critical raw material used in airframe manufacturing. The
    Materials sector in US equities is heavily weighted toward producers and processors of
    industrial metals\ldots'' \\[0.3em]
    \textbf{Validator:} ``A direct supply-side shock to a critical commodity used in aerospace
    manufacturing. The transmission channel to the US Materials sector is sound.''
  \end{block}

  \vspace{0.3em}
  \textcolor{goodgreen}{\textbf{Why this example really sells the pipeline:}}
  \begin{itemize}\itemsep0.15em
    \item \textbf{Non-trivial domain routing.} The article is about Boeing (Industrials), but
          the \emph{strong} tag goes to \textbf{Materials} via titanium --- the pipeline
          follows the \emph{economic mechanism} not the company's home sector, and independently
          across three stages (ArticleMapper, ParagraphMapper, Validator) converges on the
          same conclusion.
    \item \textbf{Quantified trigger drives the signal strength.} ``25\% of world titanium''
          is specific and large, so Materials gets 0.75 strong. The same article's aerospace
          angle drives Industrials/CAC~40/DAX but only to weak 0.35--0.45, because those
          channels are indirect.
    \item \textbf{Calibrated breadth overall.} Nine tags, one strong, zero rejected --- a
          clean consensus that this article is moderately relevant to a specific neighborhood
          of aerospace-exposed assets, with titanium as the one standout.
  \end{itemize}
\end{frame}

% ---------------------------------------------------------------
\begin{frame}[fragile]{Example 4a --- China property (article text)}
  \scriptsize
  \textbf{Headline:} \emph{China Jan-Feb Property Investment Rose 3.7\% on Year (+3 updates)}
  \hfill (DJNW, 2022-03-14, shard \texttt{b})

  \vspace{0.4em}
  \begin{block}{Full article text}
    \scriptsize
    \textbf{[0]} The downturn in China's property sector extended into early 2022 as Beijing's
    curbs on borrowing in the sector triggered uncertainty on the part of developers, while
    tight credit policies continued to pressure home buyer demand. Home sales by volume, a key
    indicator of demand, \textbf{plunged 22.1\%} in the first two months of 2022 from a year
    earlier. This compares with the 5.3\% increase for 2021 and the 143.5\% gain for the same
    period last year.\\[0.3em]
    \textbf{[1]} Real estate investment in China rose 3.7\% in the January-February period,
    slowing from the 4.4\% increase for 2021 and the 38.3\% increase in the same two-month
    period last year. \textbf{New construction starts measured by floor area declined\ldots,
    down 12.2\%} from a year earlier, widening from the 11.4\% fall for 2021.
  \end{block}

  \vspace{0.4em}
  \textbf{Pipeline output:} \quad \textbf{40 accepted, 0 rejected} \quad|\quad \textbf{5 strong,
  35 weak}

  \vspace{0.4em}
  \emph{Unlike the first two examples, one article here fans out across a large slice of the
  universe. This is the pipeline's} \textbf{breadth behavior} \emph{--- we need to check each
  tag to decide whether the reach is warranted.} The next slide bins all 40 tags by the
  quality of the transmission channel.
\end{frame}

% ---------------------------------------------------------------
\begin{frame}{Example 4b --- the 40 tags, binned by transmission quality}
  \tiny
  \setlength{\tabcolsep}{4pt}
  \renewcommand{\arraystretch}{1.15}
  \begin{tabular}{p{0.20\textwidth}p{0.73\textwidth}}
    \toprule
    \textbf{Tier 1 --- Direct} & \textbf{LME Copper (0.80), Hang Seng Index (0.75), Copper HG (0.75), LME Zinc (0.75), AUD (0.75)} \\
    \textcolor{goodgreen}{\emph{Strong}} & China is the world's largest consumer of industrial metals; construction is the primary end-use
    (wiring, plumbing, cladding, galvanized steel). HSI is dominated by Chinese property developers
    and their banks. AUD captures the commodity-FX channel --- China drives iron-ore and coal
    demand, Australia's largest exports. \\
    \midrule
    \textbf{Tier 2 --- Regional \& base-metal complex} & LME Aluminium, LME Nickel, LME Lead, Materials Sector, Kospi 200, SPI 200 Australia,
    Mini Nikkei, Mini TOPIX, CAD, 10Y/3Y Aus.\ Govt.\ Bonds \\
    \emph{Defensible} & Korea/Japan export capital goods, chemicals and steel to China; Australia's
    terms of trade hinge on iron ore and coal demand from Chinese construction. Real macro
    linkages; weak/low scores are appropriate because the channel is indirect. \\
    \midrule
    \textbf{Tier 3 --- Global growth \& risk sentiment} & Brent, WTI, Gold, DAX, Mini S\&P 500, MSCI Singapore, OMXS30 Sweden, NOK, 10Y JGB,
    2Y/5Y/10Y/30Y UST \& Ultra T-Bond \\
    \emph{Stretched} & The pipeline leans on ``China is a global growth driver'' and routes the
    shock into generic risk-off / commodity-demand channels. Coherent but exactly what the prompt
    warns against as ``generic macro atmosphere.'' Correctly scored $\leq$ 0.45. \\
    \midrule
    \textbf{Tier 4 --- Tenuous} & Soybeans, Cotton, Silver, JPY, 90-Day Aus.\ Bank Bill \\
    \textcolor{red!70!black}{\emph{Questionable}} & Multi-step chains (property $\rightarrow$
    consumer confidence $\rightarrow$ meat $\rightarrow$ soy-meal). These survive validation
    because the prompt was deliberately loosened to favor over-tagging --- a downstream
    relevance-score + source filter is the intended place to handle these. \\
    \bottomrule
  \end{tabular}
\end{frame}

% ---------------------------------------------------------------
\begin{frame}{Example 4c --- zoom + a filtering recommendation}
  \scriptsize
  \begin{block}{LME Copper --- strong, 0.80, both \ (\textcolor{goodgreen}{Tier 1})}
    \scriptsize
    \textbf{AM:} ``China is the world's largest consumer of copper, and the construction/real
    estate sector is a primary driver of industrial copper demand (wiring, plumbing, structural
    components)\ldots\ transmits directly to copper prices via a reduction in projected
    industrial demand.'' \quad \textbf{Validator:} ``\ldots a direct bearish trigger for LME
    Copper demand.''
  \end{block}

  \begin{block}{Soybeans --- weak, 0.45, paragraph\_only \ (Tier 4)}
    \scriptsize
    \textbf{Validator:} ``\ldots a sustained slump in property acts as a drag on the broader
    economy, which can lead to reduced consumer confidence and lower demand for protein
    (meat), particularly pork. Since pork consumption drives soybean-meal demand, and China is
    the largest soybean importer, this is a real but indirect channel.'' \\[0.15em]
    A five-step chain (property $\rightarrow$ confidence $\rightarrow$ meat $\rightarrow$ pork
    $\rightarrow$ soy meal) \textemdash\ but note: \textbf{China is genuinely the world's largest
    soybean importer}, and the pipeline is not hallucinating facts. Whether this kind of
    macro-chain tag is \emph{useful} for a PM is a judgment call a macro trader would have to
    make.
  \end{block}

  \begin{alertblock}{Recommendation: post-hoc filter}
    \scriptsize
    Drop any mapping that is \textbf{weak \textsc{and} score $\leq$ 0.45 \textsc{and} single-source}
    (i.e.\ only one of the two mappers fired). Intuition: a tag without cross-mapper
    corroboration \emph{and} a sub-median score is exactly the ``coherent but generic'' bucket.
    On this article, the filter shrinks the tag set from \textbf{40 $\rightarrow$ 16}: it keeps
    all 5 Tier-1 strongs, the defensible regional/base-metal cluster (Materials, Kospi, LME
    Nickel/Lead, SPI~200), and the Tier-3/4 tags where both mappers independently agreed. It
    drops the lone-mapper ``global growth'' chains (gold, US Treasuries, EAFE, Sweden, JPY,
    Singapore, $\ldots$).
  \end{alertblock}

  \textcolor{goodgreen}{\textbf{Bottom line:}} reasoning is grounded, not hallucinated; open
  question is calibration. The filter above is worth testing before tightening the validator
  prompt itself.
\end{frame}

% ---------------------------------------------------------------
\begin{frame}[fragile]{Example 5a --- the ``tags everything'' article}
  \scriptsize
  \textbf{Headline:} \emph{Dollar Firms, Ukraine War Should Send It Higher --- Market Talk}
  \hfill (DJNW, 2022-03-07, shard \texttt{a})

  \vspace{0.3em}
  \begin{block}{Full article text}
    \scriptsize
    \textbf{[0]} 1344 ET --- The dollar strengthens 0.4\% against the euro and yen and the WSJ
    Dollar Index gains 0.6\%. The Ukraine war should lead to a stronger dollar via at least
    three channels, Piper Sandler says: \textbf{safe-haven flows into US assets, a favorable
    US interest rate differential and a relatively better US economic outlook.}
    ``Normally, the dollar peaks soon after a decisive Fed policy response.'' \\[0.3em]
    \textbf{[1]} ``This crisis, however, is different because the \textbf{Fed is in a tightening
    mode given high inflation}.'' A stronger dollar tightens US financial conditions, the firm
    says, but also makes imports cheaper for the US, which could help counter strong inflation
    pressures. ``In the end, we think a stronger dollar will be another factor that could
    restrain the Fed this year.''
  \end{block}

  \vspace{0.3em}
  \textbf{Pipeline output:} \textbf{93 accepted, 0 rejected} \quad|\quad 2 strong / 91 weak
  \quad|\quad \emph{fans across every asset class}

  \vspace{0.2em}
  \centering
  \textbf{Class distribution:} commodity 25, equity index 23, gov't bond 18, US equity sector
  10, currency 9, STIR 6, volatility 2.

  \vspace{0.3em}
  \raggedright
  \textbf{The 2 strong tags are two of the three assets the article names:} Euro (0.95) and
  JPY (0.85). 3M SOFR is weak --- arguably correct since the article discusses Fed
  \emph{context}, not a direct SOFR trigger. The remaining 91 tags are weak ``global macro
  atmosphere'' inferences.
\end{frame}

% ---------------------------------------------------------------
\begin{frame}{Example 5b --- stress-test of the 4c filter}
  \scriptsize
  Applying the filter from Example 4c (\texttt{weak $\wedge$ score $\leq$ 0.45 $\wedge$
  single-source}) to this article:

  \vspace{0.5em}
  \centering
  \setlength{\tabcolsep}{6pt}
  \renewcommand{\arraystretch}{1.15}
  \begin{tabular}{lccc}
    \toprule
    \textbf{Class} & \textbf{Before} & \textbf{After filter} & \textbf{Dropped} \\
    \midrule
    commodity                & 25 & 23 &  2 \\
    equity index             & 23 & 12 & 11 \\
    government bond          & 18 & 18 &  0 \\
    US equity sector         & 10 &  8 &  2 \\
    currency                 &  9 &  9 &  0 \\
    short-term interest rate &  6 &  5 &  1 \\
    volatility               &  2 &  2 &  0 \\
    \midrule
    \textbf{TOTAL}           & \textbf{93} & \textbf{77} & \textbf{16} \\
    \bottomrule
  \end{tabular}

  \vspace{0.6em}
  \raggedright
  \textcolor{goodgreen}{\textbf{What survives the filter is still too many}} (77 tags on a
  two-paragraph Market Talk blurb). The filter is deliberately gentle --- the relevance-framed
  prompt increased cross-mapper agreement, so more tags now carry the \texttt{source=both}
  label and bypass the filter. Gov't bond, currency, volatility classes lose nothing because
  every tag in those classes is already both-source. The filter is \emph{purely post-hoc} ---
  no prompt re-engineering, no re-run needed.

  \vspace{0.3em}
  \textbf{Open question:} should a two-paragraph blurb be allowed to tag 50+ assets at all?
  A hard cap would force the pipeline to surface its top picks.
\end{frame}

% ---------------------------------------------------------------
\begin{frame}[fragile]{Example 6a --- a real failure we found, and fixed}
  \scriptsize
  \textbf{Headline:} \emph{Nymex Light Crude Oil Close --- Mar 23} (DJNW, 2022-03-23, shard \texttt{c})

  \vspace{0.3em}
  $\sim$10k chars of crude-oil contract data. No prose trigger anywhere in the body. The words
  ``crude,'' ``oil,'' and ``cocoa'' each appear \textbf{zero times}.
  \vspace{0.2em}
  \tiny
  \texttt{Source: CME Group\ldots\ Data as of: 15:54 ET\\
  Contract  Open   High   Low   Last  Chg   Prev Settle\\
  May '22\hspace{0.5em}108.85 115.40 108.38 114.58 5.31\hspace{0.5em}109.27\\
  Jun '22\hspace{0.5em}105.84\ldots\hspace{2em} Dec '27\hspace{0.5em}\phantom{00}68.20\ldots}

  \vspace{0.4em}
  \scriptsize
  \begin{columns}[T,onlytextwidth]
    \begin{column}{0.49\textwidth}
      \textbf{\textcolor{red!70!black}{BEFORE the fix}} \quad (10 tags)\\[0.2em]
      \textbf{Strong @ score = 1.0:}\\
      Cocoa, Coffee, Kansas Wheat, Live Cattle, Sugar, Chicago Wheat\\[0.3em]
      \textbf{Weak:} Cotton, Feeder Cattle, Lean Hogs, Heating Oil\\[0.3em]
      \textbf{Crude Oil: 0 tags}\\[0.3em]
      All six strong tags were hallucinated. The ParagraphMapper's reasoning for Cocoa (verbatim):
      \emph{``a direct data feed of settlement prices and trading activity for Cocoa futures
      contracts on the CME/NYBOT.''} Validator accepted every mapping.
    \end{column}
    \begin{column}{0.49\textwidth}
      \textbf{\textcolor{goodgreen}{AFTER the fix}} \quad (2 tags)\\[0.2em]
      \textbf{WTI Crude Oil}, \textbf{strong, 1.00}, paragraph\_only\\[0.1em]
      Brent Crude Oil, weak, 0.55\\[0.3em]
      \textbf{Hallucinated commodities: 0}\\[0.3em]
      Every one of the six wrong strong tags is gone. The pipeline now correctly tags the
      article's own asset (\textbf{WTI Crude Oil}) as strong, alongside a weak Brent tag for
      the adjacent benchmark. Perfect routing.
    \end{column}
  \end{columns}

  \vspace{0.2em}
  \textcolor{goodgreen}{\textbf{Companion recovery:}} \emph{ICCO Cocoa Daily Price} (2022-03-21,
  shard \texttt{c}) went 0~$\rightarrow$~\textbf{Cocoa strong 0.95}; \emph{ICE Europe White
  Sugar Close} (2022-03-21, shard \texttt{c}) 0~$\rightarrow$~1 tag. Details next slide.
\end{frame}

% ---------------------------------------------------------------
\begin{frame}{Example 6b --- root cause and the two-part fix}
  \scriptsize
  \begin{block}{What the validator was actually seeing}
    \scriptsize
    Pipeline audit: \textbf{no stage saw the headline.} \texttt{run\_summarize},
    \texttt{run\_article\_level}, \texttt{run\_paragraph\_level}, and the validator's
    \texttt{\_build\_validation\_input} all constructed the article body as paragraphs-only.
    The \texttt{headline} field was read from the article dict only when writing the output
    JSON. So when the validator was asked \emph{``is this article relevant to Cocoa?''}, it
    saw: a numerical price table, the asset label \texttt{Cocoa}, and a confident
    ParagraphMapper prior. Nothing contradicted the story, so it accepted. Six times.
  \end{block}

  \begin{block}{Fix part 1: prepend [HEADLINE] to every stage}
    \scriptsize
    Prepend \texttt{[HEADLINE] <title>} to the input of every stage, plus one neutral prompt
    sentence. No rejection instruction --- existing rules cover it. Once grounded on the
    headline, the ParagraphMapper can no longer claim a number table is Cocoa settlement
    data when the title says \emph{``Nymex Light Crude Oil Close.''}
  \end{block}

  \begin{block}{Fix part 2: drop the ``past price reporting alone'' exclusion}
    \scriptsize
    The prompt's \emph{WHAT IS NOT A SIGNAL} list had a bullet rejecting ``past price
    reporting alone,'' meant to block secondary news mentions like ``gold rose 2\% yesterday''
    --- but the mappers were also applying it to \emph{primary} exchange settlement data, so
    ICE Sugar / ICE Gasoil / ICCO Cocoa got zero tags on their own-asset articles. Removed the
    bullet; the pipeline now treats an exchange's own close report as a valid relevance
    signal for the named asset (Bryan: relevance, not alpha).
  \end{block}

  \textcolor{goodgreen}{\textbf{Combined effect:}} strong tags 525 $\rightarrow$ 727,
  both-source tags 1987 $\rightarrow$ 2428, and 34 of 57 previously zero-tag non-company
  articles now fire at least one tag.
\end{frame}

% ---------------------------------------------------------------
\section{Company-specific gate}

\begin{frame}{Does the Stage~0 gate actually work?}
  \scriptsize
  \textbf{725 filtered as company-specific (73\%), 275 kept.} The gate has to do two things:
  (a) drop pure corporate noise, and (b) \emph{not} drop articles that feature a company but
  contain real macro content.

  \vspace{0.5em}
  \begin{columns}[T,onlytextwidth]
    \begin{column}{0.49\textwidth}
      \textbf{\textcolor{goodgreen}{Correctly gated out}} (5 of 725)
      \vspace{0.1em}
      \begin{itemize}\itemsep0.1em
        \item \emph{CEO Herrman Acquires 101{,}682 Of TJX Companies Inc} --- insider buy
        \item \emph{BSE: Ajanta Pharma Ltd.\ --- Compliances Reg 39(3) Details Of Loss Of
              Certificate} --- administrative filing
        \item \emph{eMagin 4Q Loss/Shr 3c} --- microcap earnings miss, no macro mention
        \item \emph{Artiva Biotherapeutics Appoints Laura Bessen, M.D., to Its Board} ---
              board appointment
        \item \emph{HK Bourse: Announcement From SINOPEC Shanghai Petrochm} --- regulatory
              disclosure, no substance
      \end{itemize}
      \vspace{0.2em}
      Also correctly gated: Moody's municipal-debt rating changes, capital raises (e.g.\
      \emph{Phoenix Copper Seeks to Raise \$60 Mln}), single-company Russia exits (\emph{Rockwell
      Automation Suspends Business in Russia}).
    \end{column}
    \begin{column}{0.49\textwidth}
      \textbf{\textcolor{goodgreen}{Correctly kept}} (company in headline, macro inside)
      \vspace{0.1em}
      \begin{itemize}\itemsep0.1em
        \item \emph{Carnival 1Q Loss/Shr \$1.66} --- earnings note that discusses Ukraine's
              impact on fuel prices $\rightarrow$ 10 tags (Brent, WTI, Gas Oil, Energy Sector,
              Cons.\ Discretionary, FX).
        \item \emph{S\&PGR: Rising Costs Loom As Key Test For Corporate Australia} --- sector
              note citing RBA normalisation + cost inflation $\rightarrow$ SPI 200 \emph{strong},
              90-Day AU Bank Bill \emph{strong}, AUD, AU gov bonds.
        \item \emph{Boeing Co.: Suspended Parts and Maintenance Support for Russian Airlines}
              --- EU sanctions content $\rightarrow$ Industrials Sector, DAX, Euro Stoxx 50,
              VSTOXX.
        \item \emph{Small Business Growth Surpasses Pre-Pandemic Levels} --- Equifax Canada
              credit report $\rightarrow$ CAD, 10Y Canada Gov Bond, Cons.\ Discretionary,
              Energy.
        \item \emph{Wabtec and BNSF Railway Biofuels Pilot} --- renewable-diesel pilot
              $\rightarrow$ Energy, Heating Oil, Corn, Soybeans (biofuels feedstock).
      \end{itemize}
    \end{column}
  \end{columns}
\end{frame}

% ---------------------------------------------------------------
\begin{frame}[fragile]{Carnival earnings --- macro hidden inside a company headline}
  \scriptsize
  \textbf{Headline:} \emph{Carnival 1Q Loss/Shr \$1.66 \textgreater CCL} \hfill (DJNW, 2022-03-22, shard \texttt{c})

  \vspace{0.2em}
  A naive company-name filter would drop this instantly (Carnival Corp, ticker, loss-per-share).
  The pipeline keeps it because Stage~0 extracts the macro substance in the body:

  \vspace{0.3em}
  \begin{block}{Stage 0 macro\_summary (excerpt)}
    \scriptsize
    Carnival reported a US GAAP net loss of \$1.9B for Q1 2022, citing Omicron.
    \textbf{The company noted that the invasion of Ukraine is creating uncertainty and exerting
    upward pressure on fuel prices, which is materially impacting its business and liquidity.}
    Fuel consumption forecast for 2022 is 2.2M metric tons.
  \end{block}

  \vspace{0.3em}
  \textbf{Downstream tags (10 accepted, all weak):} Brent 0.45, WTI 0.40, Gas Oil 0.40, Energy
  Sector 0.35, Cons.\ Discretionary 0.55, EUR/GBP/AUD.

  \vspace{0.4em}
  \textcolor{goodgreen}{\textbf{Why this proves the gate is doing real work:}}
  \begin{itemize}\itemsep0.15em
    \item Not a keyword filter. Every surface feature screams ``drop'' --- ticker, company
          name, earnings number --- yet the article is kept because the body contains
          substantive fuel-price commentary tied to Ukraine.
    \item Downstream routing respects the economic mechanism: heaviest tags are on oil and
          fuel, \emph{not} on cruise operators. Consistent with the article's macro angle,
          not its corporate framing.
    \item All tags weak. Pipeline correctly treats earnings as a secondary source of
          fuel-price commentary, not a primary trigger.
  \end{itemize}
\end{frame}

% ---------------------------------------------------------------
\begin{frame}{Gate stress test --- macro keywords in gated headlines}
  \scriptsize
  We sampled every gated article (cs=True) whose headline contained a macro keyword (Fed, CPI,
  Ukraine, Russia, oil, gold, copper, \ldots) --- \textbf{26 articles} --- and read the bodies.

  \vspace{0.4em}
  \setlength{\tabcolsep}{4pt}
  \renewcommand{\arraystretch}{1.1}
  \begin{tabular}{p{0.33\textwidth}p{0.62\textwidth}}
    \toprule
    \textbf{Headline (gated)} & \textbf{Verdict} \\
    \midrule
    \emph{Barrick Gold Corp.\ Stock Rises Tuesday} & \textcolor{goodgreen}{\textbf{Correct.}}
    MarketWatch auto-blurb of a single stock's price move. Matches the prompt's explicit
    non-signal: \emph{``Past price reporting alone.''} 10 Barrick Gold blurbs in the sample, all gated. \\
    \emph{Rockwell Automation Suspends Business in Russia} & \textcolor{goodgreen}{\textbf{Correct.}}
    Body verbatim: \emph{``Russia and Belarus account for less than 0.5\% of total revenue.''}
    Article itself declares it immaterial. Contrast Boeing/titanium (not gated, 25\% of world supply). \\
    \emph{ITAB Shop Concept discontinues Russia ops} & \textcolor{goodgreen}{\textbf{Correct.}}
    Retail-fixture company closing a 125-employee facility. Not a bellwether. \\
    \emph{Phoenix Copper Seeks \$60 Mln for Empire Open Pit} & \textcolor{goodgreen}{\textbf{Correct.}}
    Single-company loan-note raise; \$60M for one mine is not a macro copper supply story. \\
    \emph{Manz AG: Strong demand for battery production} & \textcolor{yellow!60!black}{\textbf{Borderline.}}
    Press release about winning a Li-ion battery order. A macro strategist could argue battery
    demand is a real Cu/Ni/Li signal; the gate treats it as a single-company order. Defensible
    but the closest call in the 26. \\
    \bottomrule
  \end{tabular}

  \vspace{0.4em}
  \textcolor{goodgreen}{\textbf{Bottom line:}} 25 of 26 are clearly correct --- all match the
  prompt's explicit ``not a signal'' categories. Manz AG is a genuine gray area, defensibly
  gated under the current rules. The gate reads \emph{body content}, not headline keywords.
\end{frame}

% ---------------------------------------------------------------
\begin{frame}{External benchmark --- DJNW's own \texttt{codes} field}
  \scriptsize
  Every DJNW article ships with a \texttt{codes} dict. The \texttt{company} field lists
  tickers --- ticker-based ground truth for ``DJNW thinks this article is about a specific
  company.''

  \vspace{0.3em}
  \textbf{DJNW's own tagging:} \textbf{67.3\%} of the 1000 articles have a company ticker.
  \textbf{Our gate:} 72.6\% cs=True. \textbf{Within 5.3\,pp.}

  \vspace{0.4em}
  \centering
  \setlength{\tabcolsep}{8pt}
  \renewcommand{\arraystretch}{1.15}
  \begin{tabular}{lrr}
    \toprule
                            & \textbf{Our gate = True} & \textbf{Our gate = False} \\
    \midrule
    DJNW has company ticker & 640 & 33 \\
    DJNW no company ticker  &  86 & 241 \\
    \bottomrule
  \end{tabular}

  \vspace{0.4em}
  \raggedright
  \textcolor{goodgreen}{\textbf{Raw agreement 88.1\%, recall 95.1\%}} on DJNW-company
  articles. The 12\% disagreement decomposes cleanly:

  \vspace{0.3em}
  \textbf{33 we keep despite a company ticker} --- the intended macro-in-body cases: Apple /
  Boeing / Sylvamo / Wintershall Russia sanctions; Adidas, Leumi, Arcos Dorados, Sembcorp,
  Taiheiyo Cement earnings with macro content; \emph{Heard on the Street} columns; Wabtec
  biofuels; JD Logistics oil-price pressure; nickel surge + Chinese manufacturers.

  \textbf{86 we gate despite no ticker} --- DJNW's ticker field can't express these:
  23 closing-price bulletins, 10 Moody's/Fitch/S\&PGR municipal bond ratings, 10 exchange
  data reports, 8 small/private-company press releases, 3 WSJ culture pieces, 32 other.
  All defensibly non-macro.

  \vspace{0.2em}
  \textcolor{goodgreen}{\textbf{Takeaway:}} 5pp from DJNW's own ticker rate; disagreement
  is entirely (a) deliberate macro-body keeps and (b) non-ticker corporate content.
\end{frame}

% ---------------------------------------------------------------
\section{Aggregate diagnostics}

\begin{frame}{Relevance score distribution (all 4{,}019 accepted tags)}
  \scriptsize
  \centering
  \setlength{\tabcolsep}{4pt}
  \renewcommand{\arraystretch}{1.1}
  \begin{tabular}{lrrrrrrrrrrrrrrr}
    \toprule
    \textbf{score} & 0.30 & 0.35 & 0.40 & 0.45 & 0.50 & 0.55 & 0.60 & 0.65 & 0.70 & 0.75 & 0.80 & 0.85 & 0.90 & 0.95--1.0 \\
    \midrule
    weak   (3292) & 130 & 368 & \textbf{716} & \textbf{1108} & 318 & \textbf{494} &  55 &  68 &  29 &   3 &   . &   . &   . &   . \\
    strong ( 727) &   . &   . &   . &   . &   . &   . &   . &   . &   6 & 120 &  83 & \textbf{164} &  78 & \textbf{276} \\
    \bottomrule
  \end{tabular}

  \vspace{0.5em}
  \raggedright
  \textbf{Honest observations:}
  \begin{itemize}\itemsep0.15em
    \item \textbf{Anchoring at specific values.} Weak scores cluster at 0.40/0.45/0.55; strong at
          0.75/0.85/0.95. Earlier experiments with guard text in the prompt (``do not default
          to 0.85'') made this \emph{worse}, not better --- the mapper just anchored elsewhere.
    \item \textbf{The strong/weak boundary is clean.} Strong scores are bounded below at 0.70,
          weak bounded above at 0.75, with only minor overlap. The \emph{label} is doing more
          work than the raw score.
    \item \textbf{Implication:} treat \texttt{(signal, score)} as a pair, not as a continuous
          variable. Downstream: bin into 3--4 buckets (strong/high, strong/low, weak/med,
          weak/low) rather than use scores as a float.
  \end{itemize}
\end{frame}

% ---------------------------------------------------------------
\begin{frame}{Validator catches a factual hallucination --- 3-Month SARON}
  \scriptsize
  \textbf{Article:} \emph{Global Forex and Fixed Income Roundup: Market Talk} \hfill (DJNW,
  2022-03-02, shard \texttt{a}) --- discusses BoC rate hike and Desjardins' 2022 BoC outlook.
  ParagraphMapper proposed \textbf{3-Month SARON} via a Canadian central-bank chain.

  \vspace{0.3em}
  \begin{block}{ParagraphMapper [paragraph 7] --- the error}
    \scriptsize
    ``The transmission channel is direct: \textbf{the Bank of Canada (BOC) controls the short
    end of the Canadian yield curve}\ldots\ primary drivers for the 3-Month SARON (Swiss
    Average Rate Overnight, though in this context, \emph{the user is likely referring to the
    Canadian equivalent or the specific short-term rate benchmark for Canada given the BOC
    context})\ldots''
  \end{block}

  \begin{block}{Validator --- \textcolor{red!70!black}{REJECTED}}
    \scriptsize
    ``The asset is `3-Month SARON', the Swiss Average Rate Overnight --- \textbf{the benchmark
    interest rate for Switzerland}. The ParagraphMapper's claim in paragraph [7] that
    \textbf{the BOC controls the `3-Month SARON' is factually incorrect; the BOC controls the
    Canadian overnight rate (CORRA), not SARON}. There is no mention of the Swiss National
    Bank, Swiss inflation, or Swiss economic conditions in the article.''
  \end{block}

  \vspace{0.2em}
  \textcolor{goodgreen}{\textbf{Why this matters:}} the ParagraphMapper even flags its own
  uncertainty (``the user is likely referring to the Canadian equivalent'') but still proposes
  the mapping. The Validator --- the \emph{same} Gemma 4 26B A4B model, re-prompted --- catches
  the factual error, names the correct central bank (SNB) and the correct Canadian rate
  (CORRA), and rejects. Multi-stage self-correction as intended.
\end{frame}

% ---------------------------------------------------------------
\begin{frame}{Source $\times$ accept/reject --- what the validator actually does}
  \scriptsize
  Each (article, asset) pair reaches the validator with a \textbf{source} label recording which
  mappers fired:

  \vspace{0.5em}
  \centering
  \setlength{\tabcolsep}{8pt}
  \renewcommand{\arraystretch}{1.2}
  \begin{tabular}{lrrrr}
    \toprule
    \textbf{Source}       & \textbf{Accepted} & \textbf{Rejected} & \textbf{Total} & \textbf{Reject rate} \\
    \midrule
    both (AM $\cap$ PM agreed)     & 2{,}428 &  12 & 2{,}440 &  0.5\% \\
    article\_only (only AM fired)  &    850 &  33 &    883 &  3.7\% \\
    paragraph\_only (only PM fired)&    741 &  93 &    834 & \textbf{11.2\%} \\
    \midrule
    \textbf{total}                 & \textbf{4{,}019} & \textbf{138} & \textbf{4{,}157} & 3.3\% \\
    \bottomrule
  \end{tabular}

  \vspace{0.6em}
  \raggedright
  \textbf{Read-outs:}
  \begin{itemize}\itemsep0.15em
    \item \textbf{Cross-mapper agreement is an extremely strong prior.} When both mappers
          independently flag the same asset, the validator rejects just 0.5\% of the time
          --- the two mappers effectively pre-corroborate each other.
    \item \textbf{The ParagraphMapper is still the noisiest component.} Its lone tags get
          rejected 22$\times$ more often than two-mapper tags, which is the empirical basis
          for the ``single-source filter'' recommendation from Example 4c.
    \item \textbf{The ArticleMapper is better calibrated solo} (3.7\% reject), which makes
          sense: it sees the whole article, not one paragraph stripped of context.
  \end{itemize}
\end{frame}

% ---------------------------------------------------------------
\begin{frame}{Asset universe coverage}
  \scriptsize
  \textbf{All 95 assets are tagged at least once}, but the distribution is long-tailed.

  \vspace{0.4em}
  \centering
  \setlength{\tabcolsep}{6pt}
  \renewcommand{\arraystretch}{1.1}
  \begin{tabular}{lrrrrr}
    \toprule
    \textbf{Class} & \textbf{\# assets} & \textbf{Tags} & \textbf{Strong} & \textbf{Tags/asset} & \textbf{Strong/asset} \\
    \midrule
    commodity                & 25 & 1{,}127 & 300 & 45.1 & 12.0 \\
    equity index             & 23 &   882 & 133 & 38.3 & 5.8 \\
    government bond          & 18 &   882 & 148 & 49.0 & 8.2 \\
    currency                 &  9 &   549 &  44 & 61.0 & 4.9 \\
    US equity sector         & 11 &   350 &  49 & 31.8 & 4.5 \\
    short-term interest rate &  7 &   121 &  29 & 17.3 & 4.1 \\
    volatility               &  2 &   108 &  24 & 54.0 & 12.0 \\
    \bottomrule
  \end{tabular}

  \vspace{0.4em}
  \raggedright
  \textbf{Head (most-tagged):} CAD (103), Brent (87), WTI (81), Gold (81), AUD (78), Euro (77),
  Heating Oil (77), Energy Sector (70), Euro Bobl (69), ICE Gas Oil (68). \\[0.2em]
  \textbf{Tail (4--11 tags over 1000 articles):} 3-Month SARON (4), Health Care Sector (5),
  Cocoa (5), 90-Day NZ Bank Bill (5), Coffee (7), Cotton (10), Comm.\ Services Sector (11).

  \vspace{0.3em}
  \textbf{Observations:} energy complex and G10 FX dominate (predictable for March 2022); the
  tail is regional short rates and niche softs; STIR has lowest tags/asset because
  SOFR/SONIA/ESTR dominate; volatility 54 tags/asset is just VIX + VSTOXX.
\end{frame}

% ---------------------------------------------------------------
\section{Cost \& scale}

% ---------------------------------------------------------------
\begin{frame}{Per-stage call budget (baseline measurement)}
  \scriptsize
  \textbf{Pipeline calls for 1000 DJNW articles:}

  \vspace{0.3em}
  \centering
  \setlength{\tabcolsep}{8pt}
  \renewcommand{\arraystretch}{1.2}
  \begin{tabular}{llrr}
    \toprule
    \textbf{Stage} & \textbf{What it does} & \textbf{LLM calls} & \textbf{Wall-clock} \\
    \midrule
    Stage 0 --- Summarize & 1 call per article (macro summary, company gate) & 1{,}000 & 1:56 \\
    Stage 1 --- ArticleMapper & 1 call per (non-company article $\times$ 95 assets) & 26{,}125 & 3:53 \\
    Stage 2 --- ParagraphMapper & 1 call per (paragraph $\times$ 95 assets) & 107{,}065 & 10:26 \\
    Stage 3 --- Validator & 1 call per candidate (article, asset) pair & 3{,}484 & 2:18 \\
    \midrule
    \textbf{TOTAL}         &  & \textbf{$\sim$137{,}700} & \textbf{18:57} \\
    \bottomrule
  \end{tabular}

  \vspace{0.4em}
  \raggedright
  Throughput: $\sim$121 calls/sec on 1$\times$ B200, batched through vLLM 0.19 with prefix
  caching. \textbf{Stage 2 is 78\% of all calls} --- the dominant cost and the obvious lever
  for scaling. The next slide translates this into wall-clock estimates for the full corpus.
\end{frame}

\begin{frame}{Full-DJNW scaling --- 1996+, Bouchet fleet}
  \scriptsize
  \textbf{Corpus:} 29.11M DJNW articles (1996-01 $\rightarrow$ 2025-08); trimming the pre-1996
  head (2.56M of mostly headline-ticker feeds) removes $\sim$8\% of the total workload. Full
  corpus (1979--2025) = 31.67M for reference.

  \vspace{0.3em}
  \textbf{Baseline:} 1.137 sec/article end-to-end on 1$\times$ B200, with the 73\%
  company-specific gate already baked in. 1996+ workload is
  \textbf{$\sim$9{,}200 GPU-hours}; full corpus is $\sim$10{,}000.

  \vspace{0.3em}
  \textbf{Bouchet fleet:} \texttt{gpu\_b200} (56 B200s, $\sim$3{,}170 art/hr/GPU measured),
  \texttt{gpu\_h200} (72 H200s, $\sim$2{,}800 est.), \texttt{gpu\_rtx6000}
  (56 RTX Pro 6000 Blackwells, $\sim$2{,}400 measured). \textbf{Fleet total $\sim$184 GPUs},
  pipeline embarrassingly parallel across month shards.

  \vspace{0.3em}
  \centering
  \setlength{\tabcolsep}{8pt}
  \renewcommand{\arraystretch}{1.1}
  \begin{tabular}{lrr}
    \toprule
    \textbf{Scenario (1996+ corpus)} & \textbf{Wall clock} & \textbf{w/ Stage-2 opt.} \\
    \midrule
    1$\times$ B200 baseline          & 383 days & 191 days \\
    8 B200s (small sbatch)           &  48 days &  24 days \\
    16 B200s                         &  24 days &  12 days \\
    32 B200s (fair share)            &  12 days & \textbf{6.0 days} \\
    Full B200 partition (56)         & 6.8 days & \textbf{3.4 days} \\
    All H200s (72)                   & 6.0 days & 3.0 days \\
    \textbf{Full fleet (184 mixed)}  & \textbf{2.4 days} & \textbf{1.2 days} \\
    \bottomrule
  \end{tabular}

  \vspace{0.2em}
  \raggedright
  \textcolor{goodgreen}{\textbf{Stage 2 is the dominant cost lever.}} Two realistic levers:
  asset-universe trim ($\sim$15--20\%, pending Bryan) and a paragraph-level prefilter. The
  ``smaller Stage-2 model'' idea did \emph{not} pan out --- E4B ablation next.
\end{frame}

% ---------------------------------------------------------------
\begin{frame}{Small-model ablation program}
  \scriptsize
  Question: can a smaller model replace Gemma 4 26B A4B on Stage 2 without tanking quality?
  Same 1000 DJNW articles, seed~42, same prompts.

  \vspace{0.3em}
  \tiny
  \setlength{\tabcolsep}{3pt}
  \renewcommand{\arraystretch}{1.1}
  \begin{tabular}{p{0.15\textwidth}p{0.19\textwidth}p{0.27\textwidth}p{0.30\textwidth}}
    \toprule
    \textbf{Model} & \textbf{Architecture} & \textbf{Why tested} & \textbf{Result} \\
    \midrule
    Gemma 4 26B A4B \emph{(baseline)}
    & MoE, 4B active, server-tuned
    & \emph{baseline}
    & \textbf{18:57} wall clock, 4019 acc, 727 strong, 138 rej \\[0.15em]
    \textbf{Gemma 4 E4B}
    & dense 4B, edge-NPU (PLE, selective activation)
    & drop-in speedup attempt
    & \textcolor{red!70!black}{\textbf{failed}}: 33:33 wall, +84\% tags, 5$\times$ more strong,
      cross-currency hallucination (SARON on BoE) \\[0.15em]
    \textbf{Qwen 3.5 9B}
    & dense 9B, server-tuned
    & non-Gemma quality datapoint; not a speed candidate
    & \emph{[pending --- running on B200, gate verified 72.9\%]} \\[0.15em]
    \textbf{Qwen 3.5 4B}
    & dense 4B, server-tuned
    & true speedup candidate; isolates architecture vs param count
    & \emph{[pending --- queued on RTX]} \\
    \bottomrule
  \end{tabular}

  \vspace{0.4em}
  \scriptsize
  \raggedright
  \textbf{Why E4B failed:} 26B A4B is already at the 4B-active compute floor, so ``4B edge
  dense'' brings no speedup. Edge-NPU architectures are tuned for mobile efficiency, not
  server tensor cores. Quality also collapsed --- pipeline is reasoning-bound.

  \vspace{0.15em}
  \textbf{What Qwen 3.5 will tell us:} if 4B server-dense works $\rightarrow$ E4B failed due
  to edge tuning, not param count. If it fails $\rightarrow$ reasoning floor is above 4B
  regardless of architecture, and 26B A4B is irreducible for this task.
\end{frame}

% ---------------------------------------------------------------
\begin{frame}{DJNW input shapes --- a running taxonomy}
  \tiny
  \setlength{\tabcolsep}{4pt}
  \renewcommand{\arraystretch}{1.12}
  \begin{tabular}{p{0.22\textwidth}p{0.30\textwidth}p{0.42\textwidth}}
    \toprule
    \textbf{Shape} & \textbf{Example} & \textbf{Pipeline behavior} \\
    \midrule
    Narrative analysis column & \emph{Heard on the Street: Russia Can't Fly Without The West}
    & Multi-paragraph WSJ commentary. Best-case input: evidence localization + domain routing
      (titanium $\rightarrow$ Materials). \\
    Market Talk blurb & \emph{BOE's Rate Decision Could Benefit Sterling (HSBC)}
    & Single-paragraph analyst quote, $\sim$500--1000 chars. Handled well; AM and PM
      converge. \\
    Market news update & \emph{China Jan-Feb Property Investment Rose 3.7\%}
    & Short news story with data + commentary. The ``breadth'' case --- fans across 40+ tags;
      post-hoc filter handles the tail. \\
    Agricultural data release & \emph{USDA U.S. Export Sales and Shipments: Soybeans}
    & Structured country-by-country export table. Evidence localized to the data paragraph
      (0.95) vs the header (0.50). \\
    Exchange settlement report & \emph{ICE White Sugar Close, Nymex Crude Oil Close, Comex
    Gold Close}
    & Pure numeric tables, no prose trigger. \emph{Was} the canonical failure case
      (hallucination or silence); post-fix correctly tags the named asset (e.g.\ WTI 1.00
      strong on the Nymex article). \\
    Benchmark yield snapshot & \emph{10-Yr Benchmark Govt Yields --- U.S.\ vs Other Nations}
    & Cross-country yield table. Weak-signal data point across relevant sovereigns + FX. \\
    Market roundup / digest & \emph{Global Forex and Fixed Income Roundup: Market Talk}
    & Multi-topic, multi-region aggregation. Dangerous cross-region reasoning ---
      \textbf{where the SARON factual catch lives}. Validator is the backstop. \\
    Company earnings w/ macro & \emph{Carnival 1Q Loss/Shr \$1.66}
    & Corporate headline, macro body (Ukraine fuel pressure). Gate correctly keeps it;
      routing goes to oil/fuel complex, not just Cons.\ Discretionary. \\
    Pure corporate event & \emph{CEO Herrman Acquires 101k Of TJX}; \emph{HK Bourse
    Announcement}; \emph{BSE Compliance}
    & Insider trades, regulatory filings, MarketWatch auto-blurbs.
      \textbf{Stage-0 gate drops 73\% of the DJNW universe in this bucket.} \\
    \bottomrule
  \end{tabular}

  \vspace{0.3em}
  \raggedright
  \emph{This list is not exhaustive --- it's a running inventory of shapes we have verified
  so far. New shapes will surface as we expand to other months, other wire services, and
  eventually other text corpora (earnings transcripts, central bank speeches, research notes).}
\end{frame}

\end{document}
